% =====================================================================
\chapter{La derivada y las reglas de derivación}
\label{tema:4}
% =====================================================================

La derivada mide la tasa de cambio instantánea de una función: geométricamente, la pendiente de la recta tangente; físicamente, la velocidad. Este tema define la derivada como un límite, demuestra que la derivabilidad implica la continuidad, demuestra las reglas de derivación (suma, producto, cociente, cadena e inversa) y las aplica a las funciones elementales y a las funciones de activación. Termina con la derivación implícita y logarítmica y con las derivadas sucesivas. Se siguen los capítulos~9, 10 y~12 de Spivak~\cite{spivak}, el capítulo~4 de Apostol~\cite{apostol}, las §2.7--2.8 y el capítulo~3 de Stewart~\cite{stewart} y la unidad~1 del curso 18.01SC del MIT~\cite{mit1801}.

\section{Definición de derivada}

\begin{definicion}[Derivada]\label{def:derivada}
Sea $f$ definida en un intervalo abierto que contiene a $a$. Se dice que $f$ es \emph{derivable} en $a$ si existe y es finito el límite
\[
f'(a)=\lim_{h\to0}\frac{f(a+h)-f(a)}{h}=\lim_{x\to a}\frac{f(x)-f(a)}{x-a}.
\]
El número $f'(a)$ es la \emph{derivada} de $f$ en $a$, y la recta $y=f(a)+f'(a)(x-a)$ es la \emph{recta tangente} a la gráfica en $(a,f(a))$. La función $f'\colon x\mapsto f'(x)$, definida donde $f$ es derivable, es la \emph{función derivada}; también se escribe $\dfrac{\dd f}{\dd x}$.
\end{definicion}

El cociente $\frac{f(a+h)-f(a)}{h}$ es la pendiente de la recta secante que une $(a,f(a))$ con $(a+h,f(a+h))$; la derivada es el límite de esas pendientes cuando el segundo punto se acerca al primero. Las \emph{derivadas laterales} $f'_{+}(a)$ y $f'_{-}(a)$ se definen con los límites laterales $h\to0^{+}$ y $h\to0^{-}$; $f$ es derivable en $a$ si y solo si ambas existen y coinciden.

\begin{ejemplo}[Derivadas a partir de la definición]\label{ej:def-derivada}
(a) Para $f(x)=x^{2}$: $\dfrac{(a+h)^{2}-a^{2}}{h}=2a+h\to2a$, luego $f'(a)=2a$.
(b) Para $f(x)=1/x$ con $a\ne0$: $\dfrac{1/(a+h)-1/a}{h}=\dfrac{-1}{a(a+h)}\to-\dfrac{1}{a^{2}}$.
(c) Para $f(x)=\abs{x}$ en $a=0$: el cociente es $\abs h/h$, que vale $1$ si $h>0$ y $-1$ si $h<0$. Las derivadas laterales son $1$ y $-1$, y $\abs{x}$ no es derivable en~0.
(d) Para $\ReLU$ en~0: $f'_{+}(0)=1$ y $f'_{-}(0)=0$; no es derivable en~0. Para $x\ne0$, $\ReLU'(x)=1$ si $x>0$ y $0$ si $x<0$.
\end{ejemplo}

\begin{teorema}[Derivable implica continua]\label{teo:derivable-continua}
Si $f$ es derivable en $a$, entonces $f$ es continua en $a$.
\end{teorema}

\begin{proof}
Para $x\ne a$, $f(x)-f(a)=\dfrac{f(x)-f(a)}{x-a}\,(x-a)$. Cuando $x\to a$, el primer factor tiende a $f'(a)$ y el segundo a~0, luego $f(x)-f(a)\to f'(a)\cdot0=0$.
\end{proof}

El recíproco es falso: $\abs{x}$ y $\ReLU$ son continuas en~0 y no derivables. La siguiente reformulación de la derivada es la que mejor se generaliza a varias variables (módulo~4) y simplifica las demostraciones.

\begin{proposicion}[Derivada como aproximación lineal]\label{prop:caratheodory}
Son equivalentes: (i) $f$ es derivable en $a$ con derivada $L$; (ii) existe una función $\varphi$, continua en $a$ con $\varphi(a)=L$, tal que $f(x)=f(a)+\varphi(x)(x-a)$ para todo $x$ del dominio; (iii) $f(a+h)=f(a)+Lh+r(h)$ con $r(h)/h\to0$ cuando $h\to0$.
\end{proposicion}

\begin{proof}
(i)$\Rightarrow$(ii): se define $\varphi(x)=\frac{f(x)-f(a)}{x-a}$ si $x\ne a$ y $\varphi(a)=L$; la continuidad de $\varphi$ en $a$ es exactamente (i). (ii)$\Rightarrow$(iii): $r(h)=(\varphi(a+h)-L)h$ y $r(h)/h=\varphi(a+h)-\varphi(a)\to0$. (iii)$\Rightarrow$(i): $\frac{f(a+h)-f(a)}{h}=L+\frac{r(h)}{h}\to L$.
\end{proof}

La forma (iii) dice que, cerca de $a$, $f$ se parece a la función afín $f(a)+f'(a)(x-a)$ con un error que es pequeño \emph{incluso comparado con} $h$; esta \emph{aproximación lineal} es el primer escalón del polinomio de Taylor (\cref{tema:6}).

\section{Reglas de derivación}

\begin{teorema}[Suma, producto y cociente]\label{teo:reglas}
Si $f$ y $g$ son derivables en $a$, entonces:
\begin{enumerate}
  \item $(f+g)'(a)=f'(a)+g'(a)$ y $(cf)'(a)=cf'(a)$;
  \item (regla del producto) $(fg)'(a)=f'(a)g(a)+f(a)g'(a)$;
  \item (regla del cociente) si $g(a)\ne0$, $\Bigl(\dfrac fg\Bigr)'(a)=\dfrac{f'(a)g(a)-f(a)g'(a)}{g(a)^{2}}$.
\end{enumerate}
\end{teorema}

\begin{proof}
(i) Es el álgebra de límites aplicada a los cocientes incrementales. (ii) Sumando y restando $f(a+h)g(a)$,
\[
\frac{f(a+h)g(a+h)-f(a)g(a)}{h}=f(a+h)\,\frac{g(a+h)-g(a)}{h}+g(a)\,\frac{f(a+h)-f(a)}{h}.
\]
Cuando $h\to0$, $f(a+h)\to f(a)$ por el \cref{teo:derivable-continua}, y los cocientes tienden a $g'(a)$ y $f'(a)$. (iii) Primero, para $1/g$: como $g$ es continua en $a$ y $g(a)\ne0$, $g(a+h)\ne0$ para $h$ pequeño (\cref{prop:signo}), y
\[
\frac{1/g(a+h)-1/g(a)}{h}=-\frac{g(a+h)-g(a)}{h}\cdot\frac{1}{g(a+h)g(a)}\longrightarrow-\frac{g'(a)}{g(a)^{2}}.
\]
Después, $f/g=f\cdot(1/g)$ y la regla del producto da $\frac{f'(a)}{g(a)}-\frac{f(a)g'(a)}{g(a)^{2}}$, que es la fórmula del enunciado.
\end{proof}

\begin{proposicion}[Potencias enteras]\label{prop:potencias}
Para todo $n\in\Z$, $(x^{n})'=nx^{n-1}$ (con $x\ne0$ si $n<0$).
\end{proposicion}

\begin{proof}
Para $n=0$ y $n=1$ es inmediato. Si vale para $n\ge1$, la regla del producto da $(x^{n+1})'=(x\cdot x^{n})'=x^{n}+x\cdot nx^{n-1}=(n+1)x^{n}$; por inducción, vale para todo $n\ge0$. Para $n=-m<0$, la regla del cociente da $(1/x^{m})'=-mx^{m-1}/x^{2m}=-mx^{-m-1}$.
\end{proof}

\begin{teorema}[Regla de la cadena]\label{teo:cadena}
Si $g$ es derivable en $a$ y $f$ es derivable en $g(a)$, entonces $f\circ g$ es derivable en $a$ y
\[
(f\circ g)'(a)=f'\bigl(g(a)\bigr)\,g'(a).
\]
\end{teorema}

\begin{proof}
Por la \cref{prop:caratheodory}, existe $\psi$ continua en $b=g(a)$ con $\psi(b)=f'(b)$ y $f(y)-f(b)=\psi(y)(y-b)$ para todo $y$. Tomando $y=g(x)$,
\[
\frac{f(g(x))-f(g(a))}{x-a}=\psi\bigl(g(x)\bigr)\,\frac{g(x)-g(a)}{x-a}\qquad(x\ne a).
\]
Cuando $x\to a$, $g(x)\to g(a)$ (\cref{teo:derivable-continua}) y, por la continuidad de $\psi$ en $b$, $\psi(g(x))\to\psi(b)=f'(g(a))$; el segundo factor tiende a $g'(a)$.
\end{proof}

La demostración «ingenua», que multiplica y divide por $g(x)-g(a)$, falla si $g(x)=g(a)$ para valores de $x$ arbitrariamente próximos a $a$ (por ejemplo, $g(x)=x^{2}\sin(1/x)$ en $a=0$); la función $\psi$ evita esa división. En notación de Leibniz, si $y=f(u)$ y $u=g(x)$, la regla se escribe $\dfrac{\dd y}{\dd x}=\dfrac{\dd y}{\dd u}\dfrac{\dd u}{\dd x}$.

\begin{teorema}[Derivada de la función inversa]\label{teo:inversa}
Sea $f$ continua y estrictamente monótona en un intervalo abierto $I$, derivable en $a\in I$ con $f'(a)\ne0$. Entonces $f^{-1}$ es derivable en $b=f(a)$ y
\[
(f^{-1})'(b)=\frac{1}{f'(a)}=\frac{1}{f'\bigl(f^{-1}(b)\bigr)}.
\]
\end{teorema}

\begin{proof}
Sea $y\ne b$ en $f(I)$ y $x=f^{-1}(y)$; entonces $x\ne a$ por la inyectividad, y
\[
\frac{f^{-1}(y)-f^{-1}(b)}{y-b}=\frac{x-a}{f(x)-f(a)}=\Bigl(\frac{f(x)-f(a)}{x-a}\Bigr)^{-1}.
\]
Cuando $y\to b$, $x=f^{-1}(y)\to a$ por la continuidad de $f^{-1}$ (\cref{teo:inversa-continua}), y el cociente tiende a $1/f'(a)$.
\end{proof}

Geométricamente, la gráfica de $f^{-1}$ es la simétrica de la de $f$ respecto de $y=x$, y la simetría transforma una recta de pendiente $m$ en otra de pendiente $1/m$. Si $f'(a)=0$, la tangente a $f^{-1}$ es vertical: $f(x)=x^{3}$ tiene $f'(0)=0$ y $\sqrt[3]{y}$ no es derivable en~0.

\section{Derivadas de las funciones elementales}

\begin{proposicion}[Tabla de derivadas]\label{prop:tabla}
En los dominios indicados:
\begin{enumerate}
  \item $(e^{x})'=e^{x}$ y $(\ln x)'=1/x$ para $x>0$; $(a^{x})'=a^{x}\ln a$ para $a>0$;
  \item $(x^{r})'=rx^{r-1}$ para $x>0$ y $r\in\R$;
  \item $(\sin x)'=\cos x$, $(\cos x)'=-\sin x$, $(\tan x)'=1+\tan^{2}x=1/\cos^{2}x$;
  \item $(\arcsin x)'=1/\sqrt{1-x^{2}}$ en $(-1,1)$ y $(\arctan x)'=1/(1+x^{2})$;
  \item $(\sinh x)'=\cosh x$, $(\cosh x)'=\sinh x$ y $(\tanh x)'=1-\tanh^{2}x$.
\end{enumerate}
\end{proposicion}

\begin{proof}
(i) $\dfrac{e^{x+h}-e^{x}}{h}=e^{x}\,\dfrac{e^{h}-1}{h}\to e^{x}$ por~\eqref{eq:lim-exp}. Por el \cref{teo:inversa}, con $y=e^{x}$, $(\ln)'(y)=1/e^{x}=1/y$. Como $a^{x}=e^{x\ln a}$, la regla de la cadena da $a^{x}\ln a$.
(ii) $x^{r}=e^{r\ln x}$, luego $(x^{r})'=e^{r\ln x}\cdot r/x=rx^{r-1}$.
(iii) Por la fórmula del seno de una suma,
\[
\frac{\sin(x+h)-\sin x}{h}=\sin x\,\frac{\cos h-1}{h}+\cos x\,\frac{\sin h}{h}\longrightarrow\sin x\cdot0+\cos x\cdot1,
\]
usando la \cref{prop:senx} (de $\frac{1-\cos h}{h^{2}}\to\frac12$ se sigue $\frac{\cos h-1}{h}\to0$). El coseno es análogo, y la tangente sale de la regla del cociente.
(iv) Con $y=\sin x$, $x\in(-\pi/2,\pi/2)$: $(\arcsin)'(y)=1/\cos x=1/\sqrt{1-y^{2}}$, porque $\cos x>0$ en ese intervalo. Con $y=\tan x$: $(\arctan)'(y)=1/(1+\tan^{2}x)=1/(1+y^{2})$.
(v) Se deriva la definición: $(\sinh x)'=\frac{e^{x}+e^{-x}}{2}=\cosh x$, y análogamente $\cosh$. Por la regla del cociente y la \cref{prop:hiperbolicas}(ii), $(\tanh x)'=\frac{\cosh^{2}x-\sinh^{2}x}{\cosh^{2}x}=\frac{1}{\cosh^{2}x}=1-\tanh^{2}x$.
\end{proof}

\begin{proposicion}[Derivadas de las funciones de activación]\label{prop:derivadas-activaciones}
Para todo $x\in\R$:
\[
\sigma'(x)=\sigma(x)\bigl(1-\sigma(x)\bigr),\qquad \softplus'(x)=\sigma(x),\qquad \tanh'(x)=1-\tanh^{2}x.
\]
En particular, $\sigma'(0)=\tfrac14$, $\tanh'(0)=1$, $0<\sigma'(x)\le\tfrac14$ y $0<\tanh'(x)\le1$.
\end{proposicion}

\begin{proof}
$\sigma(x)=(1+e^{-x})^{-1}$; por la regla de la cadena, $\sigma'(x)=-(1+e^{-x})^{-2}\cdot(-e^{-x})=\dfrac{1}{1+e^{-x}}\cdot\dfrac{e^{-x}}{1+e^{-x}}$, y el segundo factor es $1-\sigma(x)$. Para la softplus, $\softplus'(x)=\dfrac{e^{x}}{1+e^{x}}=\dfrac{1}{e^{-x}+1}=\sigma(x)$. La de $\tanh$ es la \cref{prop:tabla}(v). Por último, con $s=\sigma(x)\in(0,1)$, $s(1-s)=\tfrac14-(s-\tfrac12)^{2}\in(0,\tfrac14]$; y $\tanh^{2}x\in[0,1)$.
\end{proof}

Que la derivada se exprese mediante el propio valor de la función es muy conveniente al programar la retropropagación: el valor $\sigma(x)$ calculado en la pasada hacia delante basta para obtener $\sigma'(x)$ en la pasada hacia atrás.

\section{Derivación implícita y logarítmica}

Una ecuación $F(x,y)=0$ puede definir $y$ como función de $x$ cerca de un punto sin que sea posible (o cómodo) despejarla. Si se supone que $y=y(x)$ es derivable, basta derivar la ecuación respecto de $x$ con la regla de la cadena y despejar $y'$. La justificación de que tal función $y(x)$ existe y es derivable es el teorema de la función implícita, que se estudia en varias variables.

\begin{ejemplo}[Circunferencia]\label{ej:implicita}
En $x^{2}+y^{2}=25$, derivando: $2x+2yy'=0$, luego $y'=-x/y$ si $y\ne0$. En $(3,4)$, la pendiente es $-3/4$ y la tangente $y=4-\tfrac34(x-3)$, perpendicular al radio, de pendiente $4/3$.
\end{ejemplo}

\begin{ejemplo}[Derivación logarítmica]\label{ej:logaritmica}
Para derivar $f(x)=x^{x}$ ($x>0$), que no es ni potencia ni exponencial, se toma logaritmo: $\ln f(x)=x\ln x$. Derivando, $f'(x)/f(x)=\ln x+1$, luego $f'(x)=x^{x}(1+\ln x)$. El mismo método simplifica los productos y cocientes largos: si $f=\dfrac{u_1u_2\cdots u_k}{v_1\cdots v_m}$ con factores positivos, entonces
\[
\frac{f'}{f}=\sum_{i}\frac{u_i'}{u_i}-\sum_{j}\frac{v_j'}{v_j}.
\]
\end{ejemplo}

\section{Derivadas sucesivas}

Si $f'$ es a su vez derivable, su derivada es la \emph{derivada segunda} $f''=(f')'$; en general, $f^{(n)}=(f^{(n-1)})'$. Se dice que $f$ es \emph{de clase $C^{n}$} en un intervalo si $f^{(n)}$ existe y es continua en él, y de clase $C^{\infty}$ si lo es para todo $n$. Los polinomios, $e^{x}$, $\sin$, $\cos$, $\sigma$ y $\tanh$ son de clase $C^{\infty}$; la $\ReLU$ es continua pero no de clase $C^{1}$ en~$\R$.

\begin{ejemplo}[Derivadas sucesivas]\label{ej:sucesivas}
(a) $(e^{ax})^{(n)}=a^{n}e^{ax}$. (b) $(\sin x)^{(n)}=\sin(x+n\pi/2)$, porque $\sin'x=\cos x=\sin(x+\pi/2)$. (c) $\bigl(\tfrac{1}{1-x}\bigr)^{(n)}=\tfrac{n!}{(1-x)^{n+1}}$ (por inducción). (d) $\sigma''=\bigl(\sigma(1-\sigma)\bigr)'=\sigma'(1-2\sigma)=\sigma(1-\sigma)(1-2\sigma)$, que se anula en $x=0$.
\end{ejemplo}

\begin{observacion}[Errores frecuentes]\label{obs:4-errores}
(a) $(fg)'=f'g'$ y $(f/g)'=f'/g'$ son falsas; por ejemplo, $(x\cdot x)'=2x\ne1$.
(b) Olvidar la derivada interior en la regla de la cadena: $(\sin x^{2})'=2x\cos x^{2}$, no $\cos x^{2}$.
(c) Derivar $a^{x}$ como una potencia, $xa^{x-1}$, o $x^{a}$ como una exponencial; para $x^{x}$ no sirve ninguna de las dos reglas.
(d) Concluir que una función no es derivable porque la fórmula de su derivada no está definida en un punto: $f(x)=x^{2}\sin(1/x)$, $f(0)=0$, es derivable en~0 con $f'(0)=0$ (por la definición), aunque $f'(x)$ no tenga límite en~0.
(e) En la derivación implícita, derivar $y^{2}$ como $2y$ en lugar de $2yy'$.
\end{observacion}

\begin{observacion}[Conexión con el curso]
La regla de la cadena es el fundamento de la retropropagación: una red neuronal es una composición de muchas funciones, y su derivada es el producto de las derivadas de cada capa (módulos~4, 9 y~10). Las fórmulas $\sigma'=\sigma(1-\sigma)\le\tfrac14$ y $\tanh'=1-\tanh^{2}\le1$ explican el desvanecimiento del gradiente en redes profundas con estas activaciones: un producto de muchos factores menores que~1 tiende a~0 (módulo~11). La ReLU, con derivada~1 para $x>0$, evita el problema; su falta de derivabilidad en~0 se resuelve en la práctica asignando a $\ReLU'(0)$ el valor $0$ o~$1$.
\end{observacion}

\begin{lecturas}
\item Spivak~\cite{spivak}, cap.~9 «Derivatives», cap.~10 «Differentiation» y cap.~12 «Inverse Functions».
\item Apostol~\cite{apostol}, cap.~4 «Differential Calculus».
\item Stewart~\cite{stewart}, §2.7, §2.8 y cap.~3, especialmente §3.4 «The Chain Rule», §3.5 «Implicit Differentiation», §3.6 «Derivatives of Logarithmic Functions» y §3.11 «Hyperbolic Functions».
\item MIT 18.01SC~\cite{mit1801}, unidad~1: sesiones 1, 6 y 9--20.
\end{lecturas}

\section*{Ejercicios del tema 4}
\addcontentsline{toc}{section}{Ejercicios del tema 4}

\begin{ejercicio}{1}{eje:4-grafica-derivada}
Sin calcular, describa la gráfica de la derivada de: (a) una función constante; (b) $\ReLU$; (c) $\abs{x}$; (d) $\sigma$ (¿dónde es máxima $\sigma'$?). Razone por qué la derivada de una función par es impar, con un dibujo.
\end{ejercicio}

\begin{ejercicio}{2}{eje:4-definicion}
Calcule con la definición de derivada: (a) $f'(a)$ para $f(x)=x^{3}$; (b) $f'(a)$ para $f(x)=\sqrt{x}$, $a>0$; (c) $f'(1)$ para $f(x)=\dfrac{1}{2x+1}$, y escriba la recta tangente en $x=1$.
\fuente{adaptado de MIT 18.01SC, Jerison y Mattuck~\cite{mit1801ex}, ejercicio 1C-4.}
\end{ejercicio}

\begin{ejercicio}{2}{eje:4-mecanico-1}
Derive: (a) $x^{10}+3x^{5}-\pi^{3}$; (b) $(x^{3}+x)(x^{5}+x^{2})$; (c) $\dfrac{x+2}{x^{2}-1}$; (d) $\dfrac{x^{4}+1}{x}$; (e) $(x^{2}+2)^{100}$; (f) $x^{10}(x^{2}+1)^{10}$.
\fuente{adaptado de MIT 18.01SC, Jerison y Mattuck~\cite{mit1801ex}, ejercicios 1E-1, 1E-5, 1F-1 y 1F-2.}
\end{ejercicio}

\begin{ejercicio}{2}{eje:4-mecanico-2}
Derive: (a) $e^{-x^{2}/2}$; (b) $\ln(\cos x)$; (c) $\sin^{2}(3x)$; (d) $\arctan\sqrt{x}$; (e) $\ln\bigl(x+\sqrt{x^{2}+1}\bigr)$; (f) $\sigma(wx+b)$ respecto de $x$ y respecto de $w$; (g) $\tanh(x^{2})$; (h) $\ln\sigma(x)$.
\end{ejercicio}

\begin{ejercicio}{2}{eje:4-logaritmica}
Derive mediante derivación logarítmica: (a) $x^{\sin x}$ ($x>0$); (b) $(1+1/x)^{x}$ ($x>0$); (c) $\dfrac{(x+1)^{3}\sqrt{x^{2}+4}}{(2x+3)^{5}}$ ($x>0$).
\end{ejercicio}

\begin{ejercicio}{2}{eje:4-implicita}
(a) Halle la recta tangente a la astroide $x^{2/3}+y^{2/3}=4$ en el punto $(-\sqrt{27},1)$. (b) Halle los puntos de la curva $\sin x+\sin y=\tfrac12$ en los que la tangente es horizontal.
\fuente{adaptado de MIT 18.01SC, examen~1~\cite{mit1801ex}, problema~2, y Jerison y Mattuck, ejercicio 1F-5.}
\end{ejercicio}

\begin{ejercicio}{2}{eje:4-tangentes-origen}
Halle todas las rectas tangentes a la parábola $y=1+(x-1)^{2}$ que pasan por el origen.
\fuente{adaptado de MIT 18.01SC, Jerison y Mattuck~\cite{mit1801ex}, ejercicio 1C-5.}
\end{ejercicio}

\begin{ejercicio}{2}{eje:4-pegado}
Halle las constantes $a$ y $b$ que hacen derivable en todo $\R$ la función: (a) $f(x)=x^{2}+4x+1$ si $x\ge0$ y $f(x)=ax+b$ si $x<0$; (b) la misma, con el cambio de definición en $x=1$ en lugar de en $x=0$.
\fuente{adaptado de MIT 18.01SC, Jerison y Mattuck~\cite{mit1801ex}, ejercicio 1D-6.}
\end{ejercicio}

\begin{ejercicio}{2}{eje:4-sucesivas}
(a) Calcule $f^{(5)}(x)$ para $f(x)=xe^{x}$ y conjeture y demuestre por inducción una fórmula para $f^{(n)}$. (b) Calcule $\sigma''(x)$ y $\sigma'''(0)$. (c) Calcule $\tanh''(x)$ en función de $t=\tanh x$.
\fuente{apartado (a) adaptado de MIT 18.01SC, examen~1~\cite{mit1801ex}, problema~1(c).}
\end{ejercicio}

\begin{ejercicio}{3}{eje:4-producto}
Enuncie la regla del producto y demuéstrela a partir de la definición de derivada, justificando cada paso (en particular, dónde se usa la continuidad). Use después la regla del producto y la inducción para demostrar la fórmula de Leibniz $(fg)''=f''g+2f'g'+fg''$ y, en general, $(fg)^{(n)}=\sum_{k=0}^{n}\binom nkf^{(k)}g^{(n-k)}$.
\fuente{primera parte adaptada de MIT 18.01SC, examen~1~\cite{mit1801ex}, problema~4.}
\end{ejercicio}

\begin{ejercicio}{3}{eje:4-par-impar}
Sea $f\colon\R\to\R$ derivable. Demuestre que si $f$ es par, $f'$ es impar, y que si $f$ es impar, $f'$ es par. Compruébelo con $\cosh$, $\sinh$ y $\sigma(x)-\tfrac12$.
\fuente{adaptado de MIT 18.01SC, Jerison y Mattuck~\cite{mit1801ex}, ejercicio 1F-6.}
\end{ejercicio}

\begin{ejercicio}{3}{eje:4-funcional}
Sea $f\colon\R\to\R$ tal que $f(x+y)=f(x)+f(y)+x^{2}y+xy^{2}$ para todos $x,y$, y $\lim_{x\to0}f(x)/x=1$. Halle $f(0)$, demuestre que $f$ es derivable en todo punto y calcule $f'(x)$.
\fuente{adaptado de MIT 18.01SC, examen~1~\cite{mit1801ex}, problema~6.}
\end{ejercicio}

\begin{ejercicio}{3}{eje:4-raiz-doble}
Sea $p$ un polinomio. Se dice que $a$ es una \emph{raíz doble} de $p$ si $(x-a)^{2}$ divide a $p(x)$. Demuestre que $a$ es raíz doble de $p$ si y solo si $p(a)=p'(a)=0$. Aplíquelo para hallar $c$ de modo que $x^{3}-3x+c$ tenga una raíz doble.
\fuente{adaptado de Spivak~\cite{spivak}, cap.~10, problema~24.}
\end{ejercicio}

\begin{ejercicio}{3}{eje:4-x2sen}
Sea $f(x)=x^{2}\sin(1/x)$ para $x\ne0$ y $f(0)=0$. (a) Demuestre que $f$ es derivable en todo $\R$ y calcule $f'$. (b) Demuestre que $f'$ no es continua en~0. (c) Sea $g(x)=x\sin(1/x)$, $g(0)=0$: demuestre que $g$ es continua pero no derivable en~0.
\end{ejercicio}

\begin{ejercicio}{3}{eje:4-cadena-n}
Sean $f_1,\dots,f_n$ funciones derivables de $\R$ en $\R$ y $F=f_n\circ\cdots\circ f_1$. Demuestre por inducción que, con $x_0=x$ y $x_k=f_k(x_{k-1})$,
\[
F'(x)=\prod_{k=1}^{n}f_k'(x_{k-1}).
\]
Deduzca que si $F=\sigma\circ\sigma\circ\cdots\circ\sigma$ ($n$ veces), entonces $0<F'(x)\le4^{-n}$, e interprete el resultado en términos del desvanecimiento del gradiente.
\end{ejercicio}
